% !TEX root = ../main.tex
\chapter{Principal fibre bundles}\label{ch:l19}
\providecommand{\lact}{\mathbin{\triangleright}}
\providecommand{\ract}{\mathbin{\triangleleft}}
\providecommand{\lactf}{\mathbin{\blacktriangleright}}
\providecommand{\ractf}{\mathbin{\blacktriangleleft}}
\begin{paracol}{2}

\begin{sources}
\begin{tabularx}{\linewidth}{@{}l l L@{}}
\toprule
\textbf{Video} & \textbf{Length} & \textbf{Content}\\
\midrule
\seg{19}{1}{Lecture 19: Principal fibre bundles} & 2:33:32 & Motivation (frame bundle, gauge theory, densities, spinors); left and right Lie group actions, the right action defined from a left one; equivariant maps; orbits, orbit space, stabilisers, free actions; the definition of a principal $G$-bundle; the frame bundle $LM$ as a principal $GL(d,\R)$-bundle; principal bundle maps (same group, and with a group homomorphism $\rho$; the spin frame bundle); the lemma that a principal bundle map over the identity is an isomorphism; trivial principal bundles and the theorem ``trivial $\Leftrightarrow$ there is a section'', with both directions of the proof; the frame bundle of $S^2$ is not trivial.\\
\bottomrule
\end{tabularx}

\smallskip
\textbf{Transcript} \file{materials/transcripts/C19-V01.txt} (YouTube automatic captions; the blackboard symbols for the two actions are reconstructed from the verbal description: ``arrow to the right'' for the left action, ``arrow to the left'' for the right action, a filled and an empty triangle for two different actions).
\end{sources}

\begin{reading}[Where this lecture sits in the other notes]
Not treated in the other courses (only here). Bundles were introduced in \cref{ch:l06} (bundle, bundle morphism, isomorphism, section, trivial bundle: \cref{def:mf-bundle,def:mf-morphism,def:mf-iso,def:mf-section,def:mf-trivial}); the lecturer asks to revise pull-backs, sub-bundles and restrictions (\cref{sec:mf-morph,sec:mf-pullback}). It uses Lie groups (\cref{ch:l13}), representations (\cref{ch:l17}), the exponential map (\cref{ch:l18}) and the hairy ball theorem (\cref{ch:l11}). Next: the associated bundles (\cref{ch:l20}). (From the lecture titles; not checked in detail.)
\end{reading}

% ==================================================================
\section{Motivation}\label{sec:pb-motivation}

``We have come a long way from manifolds over Lie groups, Lie algebras and back, and we are now fully prepared to come to the at least technical climax of this course'': the chapter on principal fibre bundles \ts{19}{1}{0:00} \ts{19}{1}{0:25}. What is a principal fibre bundle and why is it so important a notion? Since the technical development will be heavy at some points, it is good to know first what it is all about \ts{19}{1}{0:45}. Very roughly speaking, \emph{a principal fibre bundle is a bundle whose fibre is a Lie group}; to make this precise is exactly the point of the technical developments. This is why Lie groups are so important: they provide the fibres of fibre bundles \ts{19}{1}{1:05} \ts{19}{1}{1:26}.

Why is it so immensely important? Because principal bundles allow one to understand \emph{any} fibre bundle with a fibre $F$ on which the Lie group $G$ acts (the lecturer says ``Lie algebra'' by a slip \ts{19}{1}{2:14}); the fibre bundles on which $G$ acts and that are \emph{associated} with the principal bundle are called \textbf{associated bundles} \ts{19}{1}{1:47} \ts{19}{1}{2:35} \ts{19}{1}{2:57} (\cref{ch:l20}).

Fibre bundles are ubiquitous in physics, with a few examples \ts{19}{1}{3:17}:
\begin{itemize}
  \item In \emph{general relativity} they appear with the Lorentz group $SO(1,3)$ as the group, or, if one considers spinors, with $SL(2,\C)$, the double cover of $SO(1,3)$ \ts{19}{1}{3:38}.
  \item In \emph{Yang--Mills theory} and its generalisations (non-abelian gauge theories), the co-vector potential of ``electrodynamics'' all of a sudden becomes Lie-algebra valued; in the standard model the groups are $SU(2)$ (roughly, the electroweak theory) or $SU(3)$ (the strong interaction). These theories, and the subtleties in them, can only be properly understood by considering principal fibre bundles \ts{19}{1}{3:59} \ts{19}{1}{4:20} \ts{19}{1}{4:40}.
  \item In \emph{mathematics}: studying how the fibres of a principal bundle \emph{twist} (the bundle need not be trivial) teaches a lot about the manifold; the field of K-theory is mentioned, where one studies the bundles one can construct over a manifold $M$ \ts{19}{1}{5:01} \ts{19}{1}{5:22} \ts{19}{1}{5:43} \ts{19}{1}{6:04}.
  \item In \emph{differential geometry}: the definition of a tensor density is tricky without bundles. One looks at the principal bundle whose fibres are the general linear group of the appropriate dimension (the \emph{frame bundle}), and then the corresponding associated bundle on which $GL(d,\R)$ acts on tensor fibres with an additional determinant factor: these are the densities. The heavily used tensor densities of general relativity are thus properly understood through principal bundles and their associated bundles \ts{19}{1}{6:25} \ts{19}{1}{6:46} \ts{19}{1}{7:06}.
\end{itemize}
\begin{coursenote}[K-theory and cohomology]
Not said in class: in the lecture the cohomology groups of $M$ are said to be ``essentially defined as the various bundles one can construct over $M$''. This is a loose statement: K-theory groups are built from the \emph{vector bundles} over $M$ (isomorphism classes with a sum), not the de Rham cohomology groups of \cref{ch:l12}; the two are related (for instance through Chern classes) but not the same. \ts{19}{1}{5:22}
\end{coursenote}
``So this is lots of motivation which at the moment you do not understand, but what you do understand is: there is a Lie group, a Lie group, a Lie group \ldots{} and that is the principal fibre bundle that we want to develop'' (and a joke asking to mention the motivation in the course evaluation) \ts{19}{1}{7:26}.

% ==================================================================
\section{Lie group actions on a manifold}\label{sec:pb-actions}

We need to prepare the whole business: Lie group actions on a manifold, a topic that could have been discussed with Lie groups in \cref{ch:l13} but fits here so well that it is a preparation (section 5.1 of the lecture) \ts{19}{1}{7:46} \ts{19}{1}{8:07}.

\begin{definition}[New concepts: left action of a Lie group]\label{def:pb-left}
Let $(G,\bullet)$ be a Lie group and $M$ a smooth manifold. A smooth map
\[
  \lact\colon G\times M\to M,\qquad(g,p)\mapsto g\lact p ,
\]
(smooth, because $G\times M$ is made into a smooth manifold in the usual way) satisfying
\begin{enumerate}
  \item[(i)] $\mathbb 1\lact p=p$ for all $p\in M$ (the identity does nothing),
  \item[(ii)] $g_2\lact(g_1\lact p)=(g_2\bullet g_1)\lact p$ for all $g_1,g_2\in G$ and $p\in M$,
\end{enumerate}
is called a \textbf{left $G$-action} on $M$. \ts{19}{1}{8:29} \ts{19}{1}{8:50} \ts{19}{1}{9:11} \ts{19}{1}{9:31} \ts{19}{1}{9:51} \ts{19}{1}{10:14} \ts{19}{1}{10:35}
\end{definition}
(A map defined on a Cartesian product has its symbol put in between the two arguments, ``like the addition you know from elementary school''; in (ii) the brackets are for comfort only, there is just one way to read the expression; the order $g_2,g_1$ is kept, one can act one after the other or at once with the product \ts{19}{1}{10:14} \ts{19}{1}{10:35}.) The lecture says ``left action of the Lie algebra of the Lie group $G$'' at \ts{19}{1}{10:56} by a slip: it is the group that acts. This is ``a central notion we are going to use'' \ts{19}{1}{10:56} \ts{19}{1}{11:17}.

\begin{example}[Representations are left actions]\label{ex:pb-rep-action}
Let $G$ be a Lie group, $V$ a representation space with a linear representation $R\colon G\to GL(V)$ (\cref{ch:l17}; one needs $GL(V)$ rather than $\mathrm{End}(V)$ so that $R(g)$ is invertible). Take $M=V$ (a smooth manifold) and define $g\lact v:=R(g)v$. This is a left action. \ts{19}{1}{11:17} \ts{19}{1}{11:40} \ts{19}{1}{12:03} \ts{19}{1}{12:24} \ts{19}{1}{12:44}
\end{example}
The representations of $G$ on a vector space, as usually found in the introductory lectures of physics, are a very special case of left actions \ts{19}{1}{13:08} \ts{19}{1}{13:30}. One is almost tempted to say that left actions are ``non-linear representations'', but ``non-linear'' would suggest that it could be linear: on a general manifold there is no vector space structure, so one would not even know what linear means. The name is misleading; the idea is that one loosens the restriction \ts{19}{1}{13:30} \ts{19}{1}{13:51}.

\begin{definition}[New concepts: right action of a Lie group]\label{def:pb-right}
A smooth map $\ract\colon M\times G\to M$, $(p,g)\mapsto p\ract g$, is a \textbf{right $G$-action} if
\begin{enumerate}
  \item[(i)] $p\ract\mathbb 1=p$ for all $p\in M$,
  \item[(ii)] $(p\ract g_1)\ract g_2=p\ract(g_1\bullet g_2)$ for all $g_1,g_2\in G$ and $p\in M$.
\end{enumerate}
\ts{19}{1}{14:11} \ts{19}{1}{14:31} \ts{19}{1}{14:52} \ts{19}{1}{15:12} \ts{19}{1}{15:33} \ts{19}{1}{15:54}
\end{definition}
Whether to let a group act from the left or from the right is in principle a matter of convention. In \emph{principal} bundle theory we shall always use \emph{right} actions; for the fibre bundles on which the same group acts (the associated bundles) we shall use \emph{left} actions, and the two play nicely together \ts{19}{1}{14:11} \ts{19}{1}{14:31} \ts{19}{1}{14:52}. The point of writing the group element on the right: in (ii) the order $g_1\bullet g_2$ is the order in which one reads the acting elements, whereas for a left action the first one acting ($g_1$) stands on the right of the product, so it is convenient to write the symbol on the side of the action: ``you can always keep the order, you do not need to think about it''. It is nevertheless something really different \ts{19}{1}{16:14} \ts{19}{1}{16:35}.

\begin{proposition}[A right action from a left action]\label{prop:pb-right-from-left}
If $\lact$ is a left $G$-action on $M$, then
\[
  p\ract g:=g^{-1}\lact p
\]
is a right $G$-action on $M$. \ts{19}{1}{16:35} \ts{19}{1}{16:55} \ts{19}{1}{17:16} \ts{19}{1}{17:38}
\end{proposition}
\begin{proof}
It is smooth: the inversion of a Lie group is smooth (``one of the points where this property is important; it was implicitly important before'') \ts{19}{1}{18:40}. Property (i): $p\ract\mathbb 1=\mathbb 1^{-1}\lact p=\mathbb 1\lact p=p$ \ts{19}{1}{17:58} \ts{19}{1}{18:20}. Property (ii):
\begin{align*}
  (p\ract g_1)\ract g_2&=g_2^{-1}\lact(g_1^{-1}\lact p)\\
  &=(g_2^{-1}\bullet g_1^{-1})\lact p\\
  &=(g_1\bullet g_2)^{-1}\lact p=p\ract(g_1\bullet g_2),
\end{align*}
by the left action property and because the inverse of a product is the product of the inverses \emph{in the opposite order}. \ts{19}{1}{19:01} \ts{19}{1}{19:21} \ts{19}{1}{19:42} \ts{19}{1}{20:03} \ts{19}{1}{20:23}
\end{proof}
It is crucial that one takes the inverse: it flips the order. Without the inverse the order would not be flipped and one would not get a right action. If left and right actions were the same, this would be mere notation, but ``because of these properties a right action is a different thing'' \ts{19}{1}{20:44} \ts{19}{1}{21:04}.

\subsection{The picture to keep in mind: bases and components}\label{sec:pb-bases}
Once everything is expressed in terms of principal and associated bundles, we shall see the following recipe \ts{19}{1}{21:28}. Recall that a basis $e_1,\dots,e_{\dim M}$ of a tangent space is labelled by \emph{lower} indices, the components of a vector $X\in T_pM$ by \emph{upper} indices (``this is just a convention'') \ts{19}{1}{21:51} \ts{19}{1}{22:15} \ts{19}{1}{22:38}. If one defines a new basis by a matrix $A$, $\tilde e_m=A^n{}_me_n$, then the new components are given by precisely the inverse transformation, $\tilde X^m=(A^{-1})^m{}_nX^n$, because the linear combination of the new components with the new basis must still be the same vector: the vector is the vector, only expressed differently \ts{19}{1}{23:07} \ts{19}{1}{23:27} \ts{19}{1}{23:48}. This $A$ and $A^{-1}$ will be understood as a \emph{right} action of the general linear group $GL(d,\R)$ on the basis and a \emph{left} action of the same group on the components: ``one must be the inverse of the other'' \ts{19}{1}{24:11} \ts{19}{1}{24:32}. The bundle of bases will be the \emph{frame bundle}, the principal bundle we are going to define, and the bundle of components (say an $\R^4$ bundle) is just a bundle on which the group also acts \ts{19}{1}{24:32} \ts{19}{1}{24:53}.

This is how ``a tensor is what transforms as a tensor does'' is often defined in physics; to say precisely what it means (``tensor components are things that transform as tensor components would'') one needs principal fibre bundles \ts{19}{1}{25:14} \ts{19}{1}{25:34}. This was just a remark to show where we are going.

\subsection{Equivariant maps}\label{sec:pb-equivariant}
\begin{definition}[New concepts: $\rho$-equivariant map]\label{def:pb-equivariant}
Let $(G,\bullet)$ and $(H,\star)$ be Lie groups, $\rho\colon G\to H$ a Lie group homomorphism (smooth, and $\rho(g_1\bullet g_2)=\rho(g_1)\star\rho(g_2)$; not a representation here, but a map between groups), $M$ and $N$ smooth manifolds with a left $G$-action $\lact$ on $M$ and a left $H$-action $\lactf$ on $N$. A smooth map $f\colon M\to N$ is called \textbf{$\rho$-equivariant} if the diagram
\begin{center}
\begin{tikzpicture}[>=Stealth,node distance=3.4cm,baseline=(current bounding box.center)]
\node (GM) {$G\times M$};\node (HN) [right=of GM] {$H\times N$};
\node (M) [below=2.2cm of GM] {$M$};\node (N) [below=2.2cm of HN] {$N$};
\draw[->] (GM)--node[above,font=\small]{$\rho\times f$}(HN);
\draw[->] (M)--node[below,font=\small]{$f$}(N);
\draw[->] (GM)--node[left,font=\small]{$\lact$}(M);
\draw[->] (HN)--node[right,font=\small]{$\lactf$}(N);
\end{tikzpicture}
\end{center}
commutes, where $\rho\times f$ means that $\rho$ eats the first factor and $f$ the second; that is,
\[
  f(g\lact m)=\rho(g)\lactf f(m)\qquad\text{for all }g\in G,\ m\in M .
\]
\ts{19}{1}{25:56} \ts{19}{1}{26:18} \ts{19}{1}{26:40} \ts{19}{1}{27:00} \ts{19}{1}{27:22} \ts{19}{1}{27:43} \ts{19}{1}{28:04} \ts{19}{1}{28:28} \ts{19}{1}{28:52} \ts{19}{1}{29:14} \ts{19}{1}{29:35} \ts{19}{1}{29:56} \ts{19}{1}{30:17} \ts{19}{1}{30:38} \ts{19}{1}{30:58}
\end{definition}
Commuting means: one can first act with $g$ on $m$ by the left action on $M$ and then apply $f$, or first map $g$ by $\rho$ and $m$ by $f$ and then use the action on $N$ \ts{19}{1}{30:17} \ts{19}{1}{30:38}. So equivariant maps are the structure preserving maps between manifolds when the structure is a left action of a Lie group on each manifold, the groups being related by a homomorphism: it is again the same principle as always, structure preserving maps between objects \ts{19}{1}{31:21} \ts{19}{1}{31:42} \ts{19}{1}{32:02}. One can have $H=G$, but in physics the case $G=SO(1,3)$, $H=SL(2,\C)$ occurs: if one transforms the measurement apparatus (a Lorentz transformation, ``I rotate it or I put it into a moving system'') and the apparatus measures the spin of an object, one needs to know what the Lorentz transformation does to a spin particle, so one needs a $2:1$ homomorphism relating the two groups \ts{19}{1}{32:24} \ts{19}{1}{32:45} \ts{19}{1}{33:05} \ts{19}{1}{33:27}. (The lecture states the homomorphism in the direction $SO(1,3)\to SL(2,\C)$ here; the double covering map goes the other way, from $SL(2,\C)$ to the (orthochronous) Lorentz group: see \cref{sec:pb-pbmap}.) ``It looks excessively mathematical but it is real life, at least in theoretical physics.''

\subsection{Orbits, orbit space, stabilisers, free actions}\label{sec:pb-orbits}
\begin{definition}[New concepts: orbit]\label{def:pb-orbit}
Let $\lact\colon G\times M\to M$ be a left action. The \textbf{orbit} of a point $p\in M$ is the subset of $M$
\[
  \mathcal O_p:=\{\,q\in M\mid\exists g\in G:\ q=g\lact p\,\}
\]
of all points that can be reached from $p$ by acting with some group element. \ts{19}{1}{33:49} \ts{19}{1}{34:09} \ts{19}{1}{34:31} \ts{19}{1}{34:52}
\end{definition}
(The curly $\mathcal O$ is not a topology, it is an orbit.) \textbf{Example} \ts{19}{1}{35:16}: $M=\R^2$, $G=SO(2)$, and the left action is the vector representation,
\[
  \varphi\lact(p_1,p_2)=\begin{pmatrix}\cos\varphi&\sin\varphi\\-\sin\varphi&\cos\varphi\end{pmatrix}\begin{pmatrix}p_1\\p_2\end{pmatrix}
\]
(in quotation marks, since one has to choose coordinates to write this), the usual rotations of the plane. The orbit of a point $p\neq0$ is the circle through $p$ around the origin; the origin cannot be reached from there. The orbits are subsets of the manifold: all the circles, and also one point, the origin, which is the orbit of itself \ts{19}{1}{35:38} \ts{19}{1}{36:00} \ts{19}{1}{36:20} \ts{19}{1}{36:42} \ts{19}{1}{37:02} \ts{19}{1}{37:24}.

\begin{definition}[New concepts: orbit space]\label{def:pb-orbitspace}
For $p,q\in M$ let $p\sim q$ if there exists $g\in G$ with $q=g\lact p$ (equivalently $p=g^{-1}\lact q$): two points are equivalent if they lie on the same orbit. This is an equivalence relation, and the quotient set (\cref{def:cl-quotient}) is the \textbf{orbit space}, written
\[
  M/G:=M/\!\sim .
\]
One element of the orbit space is exactly the set of points that makes up an orbit. \ts{19}{1}{37:46} \ts{19}{1}{38:07} \ts{19}{1}{38:27} \ts{19}{1}{38:47} \ts{19}{1}{39:08} \ts{19}{1}{39:30} \ts{19}{1}{39:52}
\end{definition}
Why it is an equivalence relation: reflexive, because the identity is in the group; symmetric, because every element has an inverse (if $q=g\lact p$, then $g^{-1}\lact q=g^{-1}\lact(g\lact p)=(g^{-1}\bullet g)\lact p=\mathbb 1\lact p=p$); transitivity ``we do not do'' \ts{19}{1}{38:27} \ts{19}{1}{38:47} \ts{19}{1}{39:08}. (Not said in class: if $q=g\lact p$ and $r=h\lact q$, then $r=(h\bullet g)\lact p$.) Because the quotient is defined only using the group and the action, one writes $M/G$ and does not carry the construction around. In the example the equivalence classes are all the circles and the origin: a foliation of the space into orbits. ``This orbit space will play a very important role in principal fibre bundle theory'' \ts{19}{1}{39:52} \ts{19}{1}{40:13} \ts{19}{1}{40:34}.

\begin{definition}[New concepts: stabiliser]\label{def:pb-stabiliser}
The \textbf{stabiliser} of a point $p\in M$ is the subset of the \emph{group}
\[
  S_p:=\{\,g\in G\mid g\lact p=p\,\}
\]
of all group elements that keep the point invariant. \ts{19}{1}{40:55} \ts{19}{1}{41:16} \ts{19}{1}{41:37}
\end{definition}
Always keep in mind: the orbit is a subset of the space, the stabiliser a subset of the group \ts{19}{1}{41:37}. In the same example the stabiliser of $p\ne0$ under the rotations is just the identity, whereas the stabiliser of the origin is the entire group \ts{19}{1}{41:58} \ts{19}{1}{42:18} \ts{19}{1}{42:41}.

\begin{definition}[New concepts: free action]\label{def:pb-free}
A left (or right) action is called \textbf{free} if for \emph{every} point $p\in M$ the stabiliser consists of the identity only, $S_p=\{\mathbb 1\}$. \ts{19}{1}{43:03} \ts{19}{1}{43:25} \ts{19}{1}{43:46}
\end{definition}
So the $SO(2)$-action on $\R^2$ above is \emph{not} free (the origin is fixed by everything). A free action is one where every point is stabilised only by the identity.
\begin{answer}{Isn't $\varphi=2\pi$ a non-trivial element fixing every point?}[Free action and the circle]
A student objects that, for the rotations, one can choose another element of the group, $\varphi=2\pi$, that also leaves the point where it is. The lecturer: it is a very good example of what was said before about the exponential map \cref{ch:l18}. The rotation is $\exp(\varphi E)$ with the generator $E$; the group $SO(2)$ is compact, the group manifold is $S^1$, hence $\exp$ cannot be injective: one hits the same group element for $\varphi=0,2\pi,4\pi,\dots$. But this speaks about the \emph{parameter}, not about the group: every time one gets the identity \emph{element}. So there are no other elements than the identity that fix the point; the feeling that there must be more comes from the non-injective generation of the group elements by $\exp$ \ts{19}{1}{44:06} \ts{19}{1}{44:28} \ts{19}{1}{44:49} \ts{19}{1}{45:09} \ts{19}{1}{45:30} \ts{19}{1}{45:51} \ts{19}{1}{46:11}
\end{answer}
Can an action defined by a linear representation ever be free? No: a linear map always sends the zero vector to the zero vector, so such an action is never free \ts{19}{1}{46:33} \ts{19}{1}{46:53}. What about the translations of the plane? A translation moves every point, apart from the trivial one: so the translation action is free \ts{19}{1}{46:53} \ts{19}{1}{47:15} \ts{19}{1}{47:36}.

There are other notions, \emph{transitive} actions, \emph{effective} actions and so on; they are defined along similar lines and ``nothing to be afraid of'' if one stumbles across them, but they are not needed now: for principal bundles we only need the free actions \ts{19}{1}{47:57} \ts{19}{1}{48:18} \ts{19}{1}{48:38}.
\begin{proposition}[Orbits of a free action]\label{prop:pb-orbit-G}
If the action is free, then each orbit $\mathcal O_p$ is, as a smooth manifold, diffeomorphic to the group $G$. \ts{19}{1}{49:03} \ts{19}{1}{49:27}
\end{proposition}
The lecture only states this (``I will ask you to show that on a problem sheet''). \emph{Not said in class}: the map $G\to\mathcal O_p$, $g\mapsto g\lact p$ is surjective by definition of the orbit and injective by freeness (if $g_1\lact p=g_2\lact p$ then $g_2^{-1}\bullet g_1\in S_p=\{\mathbb 1\}$); smoothness of the inverse needs an argument and is the exercise \cref{ex:pb-orbit-G}.
\begin{example}[The punctured plane]\label{ex:pb-punctured}
Let now $M=\R^2\setminus\{0\}$ and again $SO(2)$ acting by the vector representation. This is no longer a linear representation, since $M$ is not a vector space, but still a left action (``a very trivial modification of a linear representation''). The only point with non-trivial stabiliser has been removed, so the action is \emph{free}. The orbits are the circles, and each of them is diffeomorphic to $SO(2)$, because $SO(2)$ is the Lie group built on the circle $S^1$ as underlying manifold: ``in a sense it is $SO(2)$ seen purely as a smooth manifold''. \ts{19}{1}{49:50} \ts{19}{1}{50:13} \ts{19}{1}{50:33} \ts{19}{1}{50:54} \ts{19}{1}{51:14} \ts{19}{1}{51:34} \ts{19}{1}{51:56} \ts{19}{1}{52:16}
\end{example}
This is one of the purposes of free actions: one knows immediately that the orbits are isomorphic to the group \ts{19}{1}{52:16} \ts{19}{1}{52:37}.
\begin{exercise}[Orbits of a free action]\label{ex:pb-orbit-G}
Prove \cref{prop:pb-orbit-G}. (The long version is in \cref{exam:ch19}, \cref{ex:pb-full}.)
\end{exercise}

% ==================================================================
\section{Principal fibre bundles}\label{sec:pb-def}

After these preparations the central definition \ts{19}{1}{52:58}. Whenever we talk about bundles from now on we assume everything smooth: the spaces, the projection maps, all the maps; ``if I forget to mention smooth at some point you may insert it'' \ts{19}{1}{53:24}. Recall that a bundle consists of a total space, a projection map and a base space (\cref{def:mf-bundle}).

If $E$ is equipped with a right $G$-action one says that $E$ is a \emph{right $G$-space}; this is just terminology which one does not really need, but it is convenient \ts{19}{1}{54:27} \ts{19}{1}{54:51}.
\begin{definition}[New concepts: principal $G$-bundle]\label{def:pb-principal}
Let $G$ be a Lie group. A bundle $(E,\pi,M)$ is called a \textbf{principal $G$-bundle} if
\begin{enumerate}
  \item[(i)] $E$ is a right $G$-space, i.e.\ it carries a right $G$-action $\ract\colon E\times G\to E$;
  \item[(ii)] this right $G$-action is \emph{free};
  \item[(iii)] the bundle $(E,\pi,M)$ is isomorphic, as a bundle, to the bundle $(E,\rho,E/G)$, where $E/G$ is the orbit space of the right action and $\rho\colon E\to E/G$ the canonical projection, which sends a point $\varepsilon\in E$ to its equivalence class (its orbit).
\end{enumerate}
\ts{19}{1}{53:44} \ts{19}{1}{54:07} \ts{19}{1}{55:15} \ts{19}{1}{55:37} \ts{19}{1}{55:59} \ts{19}{1}{56:20} \ts{19}{1}{56:40} \ts{19}{1}{57:02} \ts{19}{1}{57:22}
\end{definition}
Isomorphism as bundles means (\cref{def:mf-iso}): there exist diffeomorphisms $u\colon E\to E$ and $f\colon M\to E/G$ such that the diagram
\begin{center}
\begin{tikzpicture}[>=Stealth,node distance=2.6cm,baseline=(current bounding box.center)]
\node (E) {$E$};\node (E2) [right=of E] {$E$};
\node (M) [below=2cm of E] {$M$};\node (EG) [below=2cm of E2] {$E/G$};
\draw[->] (E)--node[above,font=\small]{$u$}(E2);
\draw[->] (M)--node[below,font=\small]{$f$}(EG);
\draw[->] (E)--node[left,font=\small]{$\pi$}(M);
\draw[->] (E2)--node[right,font=\small]{$\rho$}(EG);
\end{tikzpicture}
\end{center}
commutes, $\rho\circ u=f\circ\pi$. ``One always has an object that consists of three elements (total space, projection, base space); an isomorphism must map each of them, and since there is additional structure, the projection, it has to be preserved: one does not need to remember this, one can always reconstruct it by thinking about what is structure preserving'' \ts{19}{1}{58:48} \ts{19}{1}{59:08} \ts{19}{1}{59:30} \ts{19}{1}{59:51} \ts{19}{1}{1:00:12} \ts{19}{1}{1:00:33} \ts{19}{1}{1:00:53} \ts{19}{1}{1:01:14}.

Why this is the right notion: one has a manifold $E$ with a free right $G$-action; every orbit is the group itself (\cref{prop:pb-orbit-G}); $E/G$ is the space of all orbits; and the fibres of $\rho$, the preimages $\rho^{-1}([\varepsilon])$ of a point of the orbit space, are exactly the orbits, hence diffeomorphic to the group. The condition says that the bundle $(E,\pi,M)$ looks like this one \ts{19}{1}{56:40} \ts{19}{1}{57:02} \ts{19}{1}{57:22} \ts{19}{1}{57:43}.
\begin{remark}[Observation: the fibres are the group]\label{rem:pb-fibres}
Since the right action is free, the fibres $\pi^{-1}(x)$ of a principal $G$-bundle are diffeomorphic to the Lie group $G$: if $\varepsilon\in\pi^{-1}(x)$, the fibre is the orbit of $\varepsilon$. And two points of the same fibre are related by a \emph{unique} group element. This is just an observation, not part of the definition. \ts{19}{1}{57:43} \ts{19}{1}{58:06} \ts{19}{1}{58:28}
\end{remark}
In this sense a principal fibre bundle is, as announced in the beginning, a bundle whose fibre is the Lie group: it is isomorphic to a bundle which is a $G$-fibre bundle. But in that isomorphism the $G$ is not immediately visible: we only know that $G$ can act on $E$. ``This is the simplest definition one can give that does what you want'' \ts{19}{1}{1:01:36} \ts{19}{1}{1:01:57}.
\begin{coursenote}[What the definition leaves implicit]
Not said in class: (1) taking $u=\mathrm{id}_E$ in the isomorphism is the usual (and the natural) choice: then the fibres of $\pi$ are exactly the orbits; for an arbitrary diffeomorphism $u$ the diagram says only that $\pi$ and $\rho$ have fibres of the same kind. (2) For $E/G$ to be a smooth manifold and $\rho$ a smooth bundle one needs more than freeness (for instance that the action is also \emph{proper}, as it is whenever $G$ is compact); the lecture does not discuss this. (3) The standard definition adds local triviality of the bundle explicitly; it is hidden here in the requirement that the bundle be isomorphic to $E\to E/G$. \ts{19}{1}{1:01:36}
\end{coursenote}
The lecturer asks to revise, until the next lecture, the definitions of pull-backs of bundles, sub-bundles and restrictions of bundles (\cref{sec:mf-morph,sec:mf-pullback}), because they will be needed \ts{19}{1}{1:01:57} \ts{19}{1}{1:02:19}.

% ------------------------------------------------------------------
\subsection{The frame bundle}\label{sec:pb-frame}
We now consider a particular example for a principal $G$-bundle, ``of the utmost importance in differential geometry'': the \textbf{frame bundle} $LM$ \ts{19}{1}{1:02:19} \ts{19}{1}{1:02:39}.
\begin{definition}[New concepts: frame, frame bundle $LM$]\label{def:pb-frame}
Let $M$ be a smooth manifold of dimension $d$. At a point $x\in M$ let
\[
  L_xM:=\{\,(e_1,\dots,e_d)\mid e_1,\dots,e_d\text{ is a basis of }T_xM\,\}
\]
the set of ordered $d$-tuples of vectors that form a basis of the tangent space: the \textbf{frames} at $x$. (It could be called the basis bundle, but it is tradition to call it frame bundle.) The \textbf{frame bundle} is the disjoint union
\[
  LM:=\bigsqcup_{x\in M}L_xM .
\]
\ts{19}{1}{1:03:02} \ts{19}{1}{1:03:24} \ts{19}{1}{1:03:45} \ts{19}{1}{1:04:49}
\end{definition}
Each $L_xM$ is in bijection with $GL(d,\R)$: if one expresses the vectors of a basis in terms of the canonical basis of $\R^d$ one gets column vectors, which, put next to each other, form an invertible matrix \ts{19}{1}{1:04:07} \ts{19}{1}{1:04:27}. (Not said in class: this uses a chosen reference frame at $x$ and is not canonical; in particular it is not a group isomorphism, $L_xM$ has no group structure of itself.)
\begin{aside}[The disjoint union: apples and pears]
Some had asked about this kind of disjoint union \ts{19}{1}{1:05:09}. The tangent bundle was defined as the disjoint union of all the $T_xM$. Think of the spaces $T_xM$ and $T_{x'}M$ at two points as entirely different spaces: the elements in the first are apples, in the second pears; put all the apples and pears into one bag, close the eyes and pull out an element of the union: ``in my hand I have an apple. Is there any doubt where the apple comes from?'' So if one takes a vector $X$ from the tangent bundle, one always knows that there exists a unique $x$ with $X\in T_xM$, and can define the projection that assigns to $X$ the base point $x$. A construction of the type $\bigcup_xT_xM\times\{x\}$ (pairs of a point and a tangent vector) makes the point explicit, but does not solve the problem of the dimension: it is first of all just a set, and a set has no dimension; one needs to make it into a topological space and then a manifold. \ts{19}{1}{1:05:30} \ts{19}{1}{1:05:50} \ts{19}{1}{1:06:11} \ts{19}{1}{1:06:32} \ts{19}{1}{1:06:52} \ts{19}{1}{1:07:13} \ts{19}{1}{1:07:34} \ts{19}{1}{1:07:55} \ts{19}{1}{1:08:15}
\end{aside}
The manifold structure is built in complete analogy to the tangent bundle (\cref{def:tb-bundle,def:tb-chart}): $LM$ is equipped with a smooth atlas inherited from the smooth atlas of $M$. One takes an element of $LM$, expresses it in a chart-induced basis, decomposes each of its vectors, and gets a square array of numbers $e^a{}_b$ which coordinatise the element, together with the $d$ chart coordinates of the base point. The dimension of the smooth manifold $LM$ is therefore
\[
  \dim LM=\dim M+(\dim M)^2=\dim M\,(1+\dim M) .
\]
\ts{19}{1}{1:08:37} \ts{19}{1}{1:08:58} \ts{19}{1}{1:09:19} \ts{19}{1}{1:09:40} \ts{19}{1}{1:10:01} \ts{19}{1}{1:10:22}

What is still missing is the projection map. The \textbf{projection} $\pi\colon LM\to M$ is the ``apples and pears'' map: an element $(e_1,\dots,e_d)$ of $LM$ lies in $L_xM$ for a \emph{unique} $x\in M$ (the union is disjoint), and $\pi(e_1,\dots,e_d):=x$. So $(LM,\pi,M)$ is a smooth bundle; ``in order to really show this in detail you have to check a lot of things, but this is the sketch'' \ts{19}{1}{1:10:42} \ts{19}{1}{1:11:02} \ts{19}{1}{1:11:23} \ts{19}{1}{1:11:46} \ts{19}{1}{1:12:06}.

Now we want to make this into a principal $G$-bundle with $G=GL(d,\R)$: first it is a bundle (as just seen), second we establish a right $GL(d,\R)$-action on the total space \ts{19}{1}{1:12:06} \ts{19}{1}{1:12:26} \ts{19}{1}{1:12:47}. For $g\in GL(d,\R)$ with matrix entries $g^m{}_n\in\R$ ($m,n=1,\dots,d$, $\det g\ne0$ ``that is our definition of $GL$'') define
\[
  (e_1,\dots,e_d)\ract g:=\bigl(g^m{}_1e_m,\ g^m{}_2e_m,\ \dots,\ g^m{}_de_m\bigr)
\]
(summation over $m$). This is just what one does to change a basis by an invertible transformation: the new first frame vector is a linear combination of all the old frame vectors, the first collection of coefficients; the second new vector is the second collection of coefficients, and so on: ``this is precisely how you transform from one basis to another'' \ts{19}{1}{1:13:09} \ts{19}{1}{1:13:31} \ts{19}{1}{1:13:52} \ts{19}{1}{1:14:13} \ts{19}{1}{1:14:33} \ts{19}{1}{1:14:54}. (It is an action of the group on the right: not said in class, with $e'_n=g_1{}^m{}_ne_m$ one gets $(e\ract g_1)\ract g_2=(g_2{}^n{}_kg_1{}^m{}_n e_m)_k=e\ract(g_1g_2)$, the matrix product $(g_1g_2)^m{}_k=g_1{}^m{}_ng_2{}^n{}_k$.)

\textbf{Is it a free action?} A free action is one in which only the identity (the Kronecker $\delta^m{}_n$) keeps the point unchanged \ts{19}{1}{1:15:35} \ts{19}{1}{1:15:55}. For a single vector one finds non-trivial $GL$ transformations that keep it invariant, and for two vectors still; in three dimensions it gets trickier; ``but if one has an entire frame, there is nothing one can do to keep every vector of the frame invariant'': so it is free \ts{19}{1}{1:15:55} \ts{19}{1}{1:16:16} \ts{19}{1}{1:16:37}. (That is the intuitive argument; the formal proof, not said in class: if $e\ract g=e$ then $g^m{}_ne_m=e_n$ for all $n$, and since the $e_m$ are linearly independent, $g^m{}_n=\delta^m{}_n$.)

\begin{proposition}[The frame bundle is a principal $GL(d,\R)$-bundle]\label{prop:pb-frame-principal}
The bundle $(LM,\pi,M)$ with the right action above is a principal $GL(d,\R)$-bundle. \ts{19}{1}{1:16:37} \ts{19}{1}{1:16:57} \ts{19}{1}{1:17:19}
\end{proposition}
What has to be checked is that there is a bundle isomorphism $(u,f)$ between $(LM,\pi,M)$ and $(LM,\rho,LM/GL(d,\R))$, an isomorphism that in particular matches the base spaces \ts{19}{1}{1:17:40} \ts{19}{1}{1:18:02} \ts{19}{1}{1:18:22}. The intuition: the orbit space consists of orbits; take a frame at a point of the manifold and act with the whole group $GL(d,\R)$ on it: one obtains \emph{any other frame at the same point}, and only at that point (a change of basis keeps the vectors in the same $T_xM$). So the orbit of a frame is the whole $L_xM$, and if one forgets the orbit (collapses it), all that remains is the point $x$. This is precisely the identification $LM/GL(d,\R)\cong M$ that we want: ``from any one frame one can go to any other frame, and that is a free action'' \ts{19}{1}{1:18:43} \ts{19}{1}{1:19:04} \ts{19}{1}{1:19:24} \ts{19}{1}{1:19:45} \ts{19}{1}{1:20:05}. The check is to be made first intuitively (is it reasonable that it could be one?) and then written down. Another check for freeness: one must be able to reach any point of a fibre from any other point of the same fibre by some action, and by a \emph{unique} group element \ts{19}{1}{1:20:05} \ts{19}{1}{1:20:26} \ts{19}{1}{1:20:47}. ``So indeed this is a principal $GL(d,\R)$-bundle, and that is very important'' \ts{19}{1}{1:20:47} \ts{19}{1}{1:21:08}.

Why important: this is how physicists say a basis is transformed; if a basis is transformed, then also the components of a vector, a tensor or a tensor density are transformed. These objects are not elements of this bundle; they will be elements of an \emph{associated} bundle, on which there will be a \emph{left} $GL(d,\R)$-action (\cref{ch:l20}). The example cannot be taken further now, but later it will nicely become the formal setting for what has been told in the first semester of physics about tensors in classical mechanics, the inertia tensor and so on \ts{19}{1}{1:21:08} \ts{19}{1}{1:21:29} \ts{19}{1}{1:21:51}.

% ==================================================================
\section{Principal bundle maps}\label{sec:pb-pbmap}

At the beginning of the course the lecturer kept emphasising that one always studies the structure preserving maps between objects, and this remains true \ts{19}{1}{1:21:51}. Recall (\cref{def:mf-morphism}) that for two bundles $(E,\pi,M)$ and $(E',\pi',M')$ a \emph{bundle map} is a pair $(u,f)$, $u\colon E\to E'$, $f\colon M\to M'$, with $\pi'\circ u=f\circ\pi$, i.e.\ the diagram commutes \ts{19}{1}{1:22:11} \ts{19}{1}{1:22:34} \ts{19}{1}{1:22:54}. A principal bundle is certainly a bundle, but it has additional structure, the right action (``sometimes suppressed in the notation, but it should not be''), which has to be taken into account \ts{19}{1}{1:23:16} \ts{19}{1}{1:24:19} \ts{19}{1}{1:24:40}.

\begin{definition}[New concepts: principal bundle map]\label{def:pb-pbmap}
Let $(P,\pi,M)$ and $(P',\pi',M')$ be principal $G$-bundles with right actions $\ract$ and $\ract'$ (the \emph{same} group, but the action on the two total spaces may be different). A pair $(u,f)$, $u\colon P\to P'$, $f\colon M\to M'$, is a \textbf{principal bundle map} if the diagram
\begin{center}
\begin{tikzpicture}[>=Stealth,node distance=2.6cm,baseline=(current bounding box.center)]
\node (PG) {$P\times G$};\node (PG2) [right=3.4cm of PG] {$P'\times G$};
\node (P) [below=2cm of PG] {$P$};\node (P2) [below=2cm of PG2] {$P'$};
\node (M) [below=2cm of P] {$M$};\node (M2) [below=2cm of P2] {$M'$};
\draw[->] (PG)--node[above,font=\small]{$u\times\mathrm{id}_G$}(PG2);
\draw[->] (P)--node[above,font=\small]{$u$}(P2);
\draw[->] (M)--node[below,font=\small]{$f$}(M2);
\draw[->] (PG)--node[left,font=\small]{$\ract$}(P);
\draw[->] (PG2)--node[right,font=\small]{$\ract'$}(P2);
\draw[->] (P)--node[left,font=\small]{$\pi$}(M);
\draw[->] (P2)--node[right,font=\small]{$\pi'$}(M2);
\end{tikzpicture}
\end{center}
commutes: the lower square is the bundle map condition, $\pi'\circ u=f\circ\pi$; the upper square is the equivariance condition,
\[
  u(p\ract g)=u(p)\ract'g\qquad\text{for all }p\in P,\ g\in G .
\]
\ts{19}{1}{1:22:34} \ts{19}{1}{1:23:37} \ts{19}{1}{1:23:59} \ts{19}{1}{1:25:01} \ts{19}{1}{1:25:22} \ts{19}{1}{1:25:42} \ts{19}{1}{1:26:03} \ts{19}{1}{1:26:25} \ts{19}{1}{1:26:47}
\end{definition}
It is only the natural thing to require that the entire diagram commutes: it seems a subtlety, but a principal bundle map is different from a bundle map, it is a stronger notion \ts{19}{1}{1:25:22} \ts{19}{1}{1:27:10}. Questions? ``Straightforward enough'' \ts{19}{1}{1:27:10}.

\textbf{The maximal generalisation} \ts{19}{1}{1:27:34}. Take a principal $G$-bundle $(P,\pi,M)$ with right action $\ract$ and another principal $G'$-bundle $(P',\pi',M')$ with right action $\ract'$ of a possibly \emph{different} group $G'$. To talk about a map one needs an extra ingredient: a way to relate the elements of $G$ to elements of $G'$, a Lie group homomorphism $\rho\colon G\to G'$ (smooth and $\rho(g_1\bullet g_2)=\rho(g_1)\bullet'\rho(g_2)$) \ts{19}{1}{1:27:55} \ts{19}{1}{1:28:15} \ts{19}{1}{1:28:36} \ts{19}{1}{1:28:56} \ts{19}{1}{1:29:17}. Given $\rho$, $u$ and $f$, the triple is a \textbf{principal bundle map} (a bundle morphism of principal bundles) if
\[
  \pi'\circ u=f\circ\pi\qquad\text{and}\qquad u(p\ract g)=u(p)\ract'\rho(g)\ \ \text{for all }p\in P,\,g\in G ,
\]
that is, if the diagram with $u\times\rho$ in the upper row commutes. One first applies $\rho$ to turn an element of $G$ into an element of $G'$, so that it can act from the right on $P'$: ``it is totally obvious once you draw the diagrams, everything becomes child's play'' \ts{19}{1}{1:29:37} \ts{19}{1}{1:29:59} \ts{19}{1}{1:30:19} \ts{19}{1}{1:30:40} \ts{19}{1}{1:31:00} \ts{19}{1}{1:31:24}. (A bundle morphism is a bundle map: ``morph is form, a morphism is a structure preserving map, that is the game we have been playing all along'' \ts{19}{1}{1:31:24}.)

\begin{application}[General relativity: spin frame bundle]
Is this generalisation only for its own sake? There is an important application \ts{19}{1}{1:31:44} \ts{19}{1}{1:32:04}. In general relativity one often looks at the frame bundle of the space-time manifold with group $GL(d,\R)$; if there is a metric one can restrict to the orthogonal group, and for Lorentzian signature to the Lorentz group $SO(1,3)$: this is equivalent to the introduction of a metric, because then only frames related by a Lorentz transformation make up the orbits, and one has to restrict the bundle to such frames \ts{19}{1}{1:32:04} \ts{19}{1}{1:32:25} \ts{19}{1}{1:32:45} \ts{19}{1}{1:33:06}. Then one can define a second principal bundle, the \textbf{spin frame bundle}, with group $SL(2,\C)$, and a homomorphism $\rho$ between $SL(2,\C)$ and $SO(1,3)$ which is $2\!:\!1$ \ts{19}{1}{1:33:06} \ts{19}{1}{1:33:26} \ts{19}{1}{1:33:47}. A principal bundle map between them then answers: if I rotate my measurement apparatus (a Lorentz transformation) what is the corresponding spin transformation on the spinors? The details need associated bundles (\cref{ch:l20}) \ts{19}{1}{1:34:08} \ts{19}{1}{1:34:28}.

Why spinors need this: there is no (finite-dimensional, unitary, irreducible) spin representation of $GL(4,\R)$, so one has to restrict to $SO(1,3)$ (or its cover), and for that one first needs the metric; tensor fields do not need a metric, they exist because the manifold is smooth, but finite-dimensional spinor representations require a geometry before one can even say what a spinor is. The principal bundle map then tells physically that, once one transforms a local frame in the frame bundle, one has to use the corresponding transformation in the spin frame bundle \ts{19}{1}{1:34:28} \ts{19}{1}{1:34:49} \ts{19}{1}{1:35:10} \ts{19}{1}{1:35:30} \ts{19}{1}{1:35:51}. Should these remarks remain mysterious: ``have patience''.
\end{application}
\begin{coursenote}[Direction of $\rho$ in the example]
By the definition just given $\rho$ goes from the group of the source bundle to the group of the target bundle, $\rho\colon G\to G'$ with $u(p\ract g)=u(p)\ract'\rho(g)$. The covering homomorphism goes from $SL(2,\C)$ to the Lorentz group, so in the application the \emph{spin frame bundle} (group $SL(2,\C)$) is the source $P$ and the (Lorentz) frame bundle the target $P'$; the lecture's labels are not consistent on this point [?]. \ts{19}{1}{1:33:47}
\end{coursenote}
\begin{answer}{Does $u$ determine $\rho$?}[The role of $\rho$]
A student asks (inaudible) about the connection between $\rho$ and $u$. The lecturer: $\rho$ is an extra piece of data. The connection is only established by the requirement that $u$ and $\rho$ be chosen so that they satisfy the condition. One cannot say that, given $u$, the $\rho$ is determined: there might be several $\rho$ making $u$ into a principal bundle map of this type; for any Lie group homomorphism one can impose the condition. The data are all prescribed and the condition is imposed, as in the example where one already knows how $SL(2,\C)$ is mapped into $SO(1,3)$; in particular circumstances $\rho$ may be determined by $u$, but in full generality not \ts{19}{1}{1:36:20} \ts{19}{1}{1:36:43} \ts{19}{1}{1:37:03} \ts{19}{1}{1:37:24} \ts{19}{1}{1:37:44} \ts{19}{1}{1:38:06}
\end{answer}

\subsection{A principal bundle map over the identity is an isomorphism}\label{sec:pb-lemma}
Principal bundle maps are a stronger notion than bundle maps, and a consequence of this is the following lemma \ts{19}{1}{1:38:06} \ts{19}{1}{1:38:28}.
\begin{lemma}[Principal bundle maps over the identity]\label{lem:pb-iso}
Let $(P,\pi,M)$ and $(P',\pi',M)$ be principal $G$-bundles over the \emph{same} base manifold $M$, and $u\colon P\to P'$ a principal bundle map with $\pi'\circ u=\pi$ (base map the identity) and $u(p\ract g)=u(p)\ract'g$. Then $u$ is bijective (and a diffeomorphism): it is a bundle isomorphism. \ts{19}{1}{1:38:49} \ts{19}{1}{1:39:09} \ts{19}{1}{1:39:30} \ts{19}{1}{1:39:52} \ts{19}{1}{1:40:12}
\end{lemma}
\begin{proof}
If one has a map from $M$ to $M$ that commutes with the projections, it is the identity (``we go to the same point''), which is why no $f$ appears \ts{19}{1}{1:42:38}. We prove injectivity and surjectivity; then $u$ is invertible, and the rest (smoothness) follows \ts{19}{1}{1:40:32} \ts{19}{1}{1:40:54}.\par
\emph{Injective.} Let $u(p_1)=u(p_2)$. Then $\pi(p_1)=\pi'(u(p_1))=\pi'(u(p_2))=\pi(p_2)$, so $p_1,p_2$ lie in the same fibre; hence there is a \emph{unique} $g$ with $p_2=p_1\ract g$ (freeness, and principal bundle: the fibre is an orbit). Then
\[
  u(p_1)=u(p_2)=u(p_1\ract g)=u(p_1)\ract'g ,
\]
by equivariance. So $g$ fixes the point $u(p_1)$, and since the action is free, $g=\mathbb 1$. Hence $p_2=p_1$. (Both the bundle map property and the extra equivariance property were used.) \ts{19}{1}{1:41:15} \ts{19}{1}{1:41:35} \ts{19}{1}{1:43:18} \ts{19}{1}{1:43:42} \ts{19}{1}{1:44:03} \ts{19}{1}{1:44:24} \ts{19}{1}{1:44:45} \ts{19}{1}{1:45:06} \ts{19}{1}{1:45:26} \ts{19}{1}{1:45:47} \ts{19}{1}{1:46:07} \ts{19}{1}{1:46:28} \ts{19}{1}{1:46:49} \ts{19}{1}{1:47:09} \ts{19}{1}{1:47:30}\par
\emph{Surjective.} Let $p'\in P'$. Choose some $p\in P$ in the fibre $\pi^{-1}(\pi'(p'))$ over the same base point; it exists because $P$ is a bundle over $M$ with non-empty fibres. Then $\pi'(u(p))=\pi(p)=\pi'(p')$ (bundle map aspect only), so $u(p)$ and $p'$ lie in the same fibre of $P'$ and there is a unique $g\in G$ with $u(p)\ract'g=p'$. By equivariance $u(p\ract g)=u(p)\ract'g=p'$, so $\tilde p:=p\ract g$ is the point we were looking for. \ts{19}{1}{1:47:52} \ts{19}{1}{1:48:13} \ts{19}{1}{1:48:33} \ts{19}{1}{1:48:57} \ts{19}{1}{1:49:17} \ts{19}{1}{1:49:38} \ts{19}{1}{1:49:59} \ts{19}{1}{1:50:19} \ts{19}{1}{1:50:40} \ts{19}{1}{1:51:00} \ts{19}{1}{1:51:21} \ts{19}{1}{1:51:44} \ts{19}{1}{1:52:06} \ts{19}{1}{1:52:27} \ts{19}{1}{1:52:47} \ts{19}{1}{1:53:09} \ts{19}{1}{1:53:29}
\end{proof}
Without the extra principal-bundle condition beyond the pure bundle map, neither injectivity nor surjectivity could be proven. ``It is really a strengthening of what a bundle map is on a principal bundle: a principal bundle map is more than a bundle map, and it makes this automatically a diffeomorphism'' \ts{19}{1}{1:53:50} \ts{19}{1}{1:54:12}.
\begin{coursenote}[Smoothness of the inverse]
Not said in class: the lecture proves bijectivity only and says that the smoothness ``also will be the case''. One way: over a small open set $U\subseteq M$ choose local sections $\sigma$ of $P$ and $\sigma'$ of $P'$ (they exist by local triviality) and write $u(\sigma(x))=\sigma'(x)\ract'\gamma(x)$ with a smooth $\gamma\colon U\to G$; then $u$ acts as $\sigma(x)\ract g\mapsto\sigma'(x)\ract'(\gamma(x)\bullet g)$, whose inverse is smooth since $G$ is a Lie group. \ts{19}{1}{1:54:12}
\end{coursenote}
In the proof we used several times the fact that two points in the same fibre of a principal bundle differ by a unique group element (\cref{rem:pb-fibres}). The lemma is the basis of the next definition and theorem \ts{19}{1}{1:54:32} \ts{19}{1}{1:54:55}.

% ==================================================================
\section{Trivial principal bundles and sections}\label{sec:pb-trivial}

\begin{definition}[New concepts: trivial principal bundle]\label{def:pb-trivial}
The \textbf{trivial principal $G$-bundle} over $M$ is the bundle $(M\times G,\pi_1,M)$, with $\pi_1(x,h):=x$ the projection onto the first factor, and the right action
\[
  (x,h)\ractf g:=(x,h\bullet g)
\]
(the base point untouched; one acts from the right by $g$ on the group factor). It is a principal $G$-bundle: the action is free because in the group there is always the inverse ($h\bullet g=h$ forces $g=\mathbb 1$). A principal $G$-bundle $(P,\pi,M)$ is called \textbf{trivial} if there exists a principal bundle map $u\colon P\to M\times G$ over $\mathrm{id}_M$ (i.e.\ $\pi_1\circ u=\pi$ and $u(p\ract g)=u(p)\ractf g$). \ts{19}{1}{1:54:55} \ts{19}{1}{1:55:24} \ts{19}{1}{1:55:44} \ts{19}{1}{1:56:05} \ts{19}{1}{1:56:25} \ts{19}{1}{1:56:45} \ts{19}{1}{1:57:06} \ts{19}{1}{1:57:26} \ts{19}{1}{1:57:47}
\end{definition}
Since a principal bundle map is automatically a diffeomorphism (\cref{lem:pb-iso}), ``trivial'' can be phrased either as: diffeomorphic as a principal $G$-bundle to $M\times G$, or as: there exists a principal bundle map to $M\times G$ \ts{19}{1}{1:55:24} \ts{19}{1}{1:55:44} \ts{19}{1}{1:59:52} \ts{19}{1}{2:00:13}. One should insist that it is \emph{not} enough to have a bundle map to the trivial bundle (one would be tempted to define it like that): one only calls it trivial if there exists a \emph{principal} bundle map, which is the stronger notion \ts{19}{1}{2:00:13} \ts{19}{1}{2:00:33}.

At the beginning of the lecture we said that a principal $G$-bundle very roughly is a bundle with fibre $G$. This trivial bundle is the very trivial way to make $G$ a fibre: one has a manifold and attaches the group at every point. A principal $G$-bundle is also, very precisely, a bundle with fibre $G$ (because of the isomorphism with $P\to P/G$), but it is certainly not always trivial like this \ts{19}{1}{1:57:47} \ts{19}{1}{1:58:07} \ts{19}{1}{1:58:28} \ts{19}{1}{1:58:48} \ts{19}{1}{1:59:12} \ts{19}{1}{1:59:32}.

Why are we interested in trivial principal bundles? One can think of them globally as having the group just attached as a fibre: very simple. The question is how common that is, and what it has to do with what we did before \ts{19}{1}{2:00:55} \ts{19}{1}{2:01:15}.
\begin{answer}{What exactly is called trivial here?}[Trivial bundle and trivial principal bundle]
A question (partly inaudible) about the colours of the diagram. The lecturer: the right-hand side $M\times G$ with the projection $\pi_1$ is certainly what one calls a trivial \emph{bundle}, because it is a Cartesian product. Any principal bundle that is diffeomorphic to it \emph{as a principal bundle}, i.e.\ by a diffeomorphism of the total spaces that preserves all the group actions, is called trivial: all the structure may look more complicated, but after an appropriate diffeomorphism it takes this simple form. It was the same with bundles before: a bundle is trivial if the total space is the Cartesian product of the base with a fibre. In the extended picture, the principal bundle $P\to M$ is, by definition, bundle-isomorphic to $P\to P/G$ with its free $G$-action, so a $G$-fibre bundle with the same fibre at every point; hence it makes sense to ask the question. Not every principal bundle is trivial; otherwise the notion would not be worth a definition \ts{19}{1}{2:01:36} \ts{19}{1}{2:01:57} \ts{19}{1}{2:02:17} \ts{19}{1}{2:02:37} \ts{19}{1}{2:02:58} \ts{19}{1}{2:03:19} \ts{19}{1}{2:03:39} \ts{19}{1}{2:04:02} \ts{19}{1}{2:04:27} \ts{19}{1}{2:04:47} \ts{19}{1}{2:05:09} \ts{19}{1}{2:05:29} \ts{19}{1}{2:05:50}
\end{answer}

\begin{theorem}[Triviality and sections]\label{thm:pb-trivial-section}
A principal $G$-bundle $(P,\pi,M)$ is trivial if and only if it admits a smooth section, i.e.\ a smooth map $\sigma\colon M\to P$ with $\pi\circ\sigma=\mathrm{id}_M$ (\cref{def:mf-section}). \ts{19}{1}{2:06:12} \ts{19}{1}{2:06:32} \ts{19}{1}{2:06:52} \ts{19}{1}{2:07:15} \ts{19}{1}{2:07:36}
\end{theorem}
The point of the theorem: finding a diffeomorphism or a principal bundle map can be tricky, and it is hard to show that it does not exist; so it is good to have a criterion, here a necessary and sufficient one \ts{19}{1}{2:06:12} \ts{19}{1}{2:06:32}.
\begin{example}[The frame bundle of the 2-sphere is not trivial]\label{ex:pb-S2}
Let $M=S^2$ and $P=LS^2$, the frame bundle, a principal $GL(2,\R)$-bundle (\cref{prop:pb-frame-principal}). A section of the frame bundle would be a frame at every point, that is a basis of the tangent space smoothly depending on the point: two linearly independent, everywhere non-vanishing smooth vector fields. But there is not even one everywhere non-vanishing smooth vector field on $S^2$: ``you cannot comb a sphere'' (the hairy ball theorem of \cref{ch:l11}). Hence there is no smooth section, and by \cref{thm:pb-trivial-section} the principal $GL(2,\R)$-bundle $LS^2$ is not trivial. So there are principal bundles that are non-trivial, and others, like $M\times G$, that are trivial. \ts{19}{1}{2:07:58} \ts{19}{1}{2:08:19} \ts{19}{1}{2:08:40} \ts{19}{1}{2:09:01} \ts{19}{1}{2:09:24} \ts{19}{1}{2:09:48} \ts{19}{1}{2:10:08} \ts{19}{1}{2:10:29} \ts{19}{1}{2:10:49}
\end{example}
(At the end of the lecture the lecturer relates this to the previous lectures: there is no basis of the module of vector fields, the $C^\infty(M)$-module of vector fields is not free, only projective, the theorem of Serre and Swan (\cref{ch:l11}): ``all these things belong together'' \ts{19}{1}{2:31:13} \ts{19}{1}{2:31:34} \ts{19}{1}{2:31:55}.) Now we want to understand the theorem better by proving it \ts{19}{1}{2:10:49} \ts{19}{1}{2:11:10}.
\begin{proof}
\emph{($\Rightarrow$)} Suppose the bundle is trivial: there is a principal bundle map $u\colon P\to M\times G$ over the identity, a diffeomorphism by \cref{lem:pb-iso}. Define
\[
  \sigma(x):=u^{-1}(x,\mathbb 1) ,
\]
that is, put the identity of the group in the second slot and push it back. It is smooth, and a section because $\pi\bigl(\sigma(x)\bigr)=\pi_1\bigl(u(\sigma(x))\bigr)=\pi_1(x,\mathbb 1)=x$, since $u$ is a bundle map. \ts{19}{1}{2:11:32} \ts{19}{1}{2:11:53} \ts{19}{1}{2:12:13} \ts{19}{1}{2:12:34} \ts{19}{1}{2:12:54}\par
\emph{($\Leftarrow$)} Suppose a smooth section $\sigma$ is given \ts{19}{1}{2:13:16} \ts{19}{1}{2:13:36}. Let $p\in P$ be any point; it lies in some fibre. The point $\sigma(\pi(p))$ lies in the \emph{same} fibre, so there is a unique group element $\chi_\sigma(p)\in G$ with
\[
  \sigma(\pi(p))\ract\chi_\sigma(p)=p .
\]
It exists, and is unique, because the action is free and the fibre is an orbit. It depends on the section $\sigma$ (another section would need another element to bridge the gap) and on the point $p$, so it is a map $\chi_\sigma\colon P\to G$ \ts{19}{1}{2:13:57} \ts{19}{1}{2:14:20} \ts{19}{1}{2:14:40} \ts{19}{1}{2:15:01} \ts{19}{1}{2:15:22} \ts{19}{1}{2:15:42} \ts{19}{1}{2:16:03} \ts{19}{1}{2:16:24} \ts{19}{1}{2:16:45} \ts{19}{1}{2:17:06}.\par
\emph{Claim:} $\chi_\sigma(p\ract g)=\chi_\sigma(p)\bullet g$ for all $g\in G$. Indeed, $\pi(p\ract g)=\pi(p)$ (the action stays in the fibre), hence $\sigma(\pi(p\ract g))=\sigma(\pi(p))$ and, by definition of $\chi_\sigma$, $\sigma(\pi(p))\ract\chi_\sigma(p\ract g)=p\ract g$. On the other hand, acting on the defining equation with $g$,
\[
  \sigma(\pi(p))\ract\bigl(\chi_\sigma(p)\bullet g\bigr)=\bigl(\sigma(\pi(p))\ract\chi_\sigma(p)\bigr)\ract g=p\ract g ,
\]
because it is a right action. So the same point is acted upon by the two elements $\chi_\sigma(p\ract g)$ and $\chi_\sigma(p)\bullet g$ and the results coincide; since the action is free, the two group elements are equal. \ts{19}{1}{2:17:27} \ts{19}{1}{2:17:48} \ts{19}{1}{2:18:10} \ts{19}{1}{2:18:32} \ts{19}{1}{2:18:53} \ts{19}{1}{2:20:00} \ts{19}{1}{2:20:22} \ts{19}{1}{2:21:03} \ts{19}{1}{2:21:24} \ts{19}{1}{2:22:09} \ts{19}{1}{2:22:30} \ts{19}{1}{2:22:52} \ts{19}{1}{2:23:13} \ts{19}{1}{2:23:33} \ts{19}{1}{2:23:53} \ts{19}{1}{2:24:14}\par
Now define
\[
  u_\sigma\colon P\to M\times G,\qquad u_\sigma(p):=\bigl(\pi(p),\chi_\sigma(p)\bigr).
\]
It is a bundle map: $\pi_1\circ u_\sigma=\pi$, obvious since $\pi_1$ projects out the first component. It is equivariant:
\[
  u_\sigma(p\ract g)=\bigl(\pi(p),\chi_\sigma(p)\bullet g\bigr)=u_\sigma(p)\ractf g ,
\]
where we used the claim (and $\pi(p\ract g)=\pi(p)$): ``this is why we had to show that if one first acts on $p$ and then uses $\chi_\sigma$, one can also first use $\chi_\sigma$ and then blob with $g$''. So $u_\sigma$ is a principal bundle map into the trivial bundle: the bundle is trivial \ts{19}{1}{2:24:36} \ts{19}{1}{2:24:56} \ts{19}{1}{2:25:17} \ts{19}{1}{2:25:37} \ts{19}{1}{2:25:58} \ts{19}{1}{2:26:20} \ts{19}{1}{2:26:41} \ts{19}{1}{2:27:01} \ts{19}{1}{2:27:22} \ts{19}{1}{2:27:43} \ts{19}{1}{2:28:04} \ts{19}{1}{2:28:25} \ts{19}{1}{2:28:46} \ts{19}{1}{2:29:07} \ts{19}{1}{2:29:28} \ts{19}{1}{2:29:48} \ts{19}{1}{2:30:11} \ts{19}{1}{2:30:32} \ts{19}{1}{2:30:53}.
\end{proof}
So: having a section means that the bundle is trivial, and conversely, if the bundle is trivial, one can construct a section. The lecture does not discuss the smoothness of $\chi_\sigma$.
\begin{coursenote}[Smoothness of $\chi_\sigma$]
Not said in class: $\chi_\sigma$ is smooth; locally one has a local section $\sigma'$ of the principal bundle with $p=\sigma'(\pi(p))\ract\chi'(p)$ for a smooth $\chi'\colon\pi^{-1}(U)\to G$, and $\sigma(x)=\sigma'(x)\ract\gamma(x)$ with $\gamma$ smooth, so $\chi_\sigma(p)=\gamma(\pi(p))^{-1}\bullet\chi'(p)$. The trial in the lecture of writing the equation for $\chi_\sigma(p\ract g)$ had a few false starts at the board (``hang on, can that be the thing I want to do?''), cleaned up here. \ts{19}{1}{2:19:13} \ts{19}{1}{2:19:34}
\end{coursenote}
\begin{aside}[Looking ahead]
The lecture ends here: ``see you after Christmas, and we will continue with associated bundles, and then you will finally understand why we are so interested in principal bundles, because they have associated bundles, and in the associated bundles live the objects we like so much: spinors, tensors, tensor densities''. If in physics we say we rotate something, we do an $SO(3)$ transformation, how do the different objects, which are sections of different associated bundles, transform under the same transformation? Understood in terms of the principal bundle: ``in this sense the principal bundle really thrones over all of the others. If you understand the principal bundle, you understand all the associated bundles'' \ts{19}{1}{2:32:17} \ts{19}{1}{2:32:38} \ts{19}{1}{2:32:58} \ts{19}{1}{2:33:19}
\end{aside}

% ==================================================================
\section*{Key concepts}
\begin{keyconcepts}
\textbf{Retrieval questions.}
\begin{enumerate}
  \item Define a left and a right $G$-action. How does one get a right action from a left one, and why is the inverse needed?
  \item What is a $\rho$-equivariant map? Write the commuting diagram and the equation.
  \item Define orbit, orbit space $M/G$ and stabiliser. Which is a subset of the space, which of the group? When is an action free?
  \item Why is the rotation action on $\R^2$ not free, and why is it free on $\R^2\setminus\{0\}$? Why is a linear action never free?
  \item State the three conditions of a principal $G$-bundle. Why are its fibres diffeomorphic to $G$?
  \item Define the frame bundle $LM$, its dimension, the projection and the right $GL(d,\R)$-action. Why is the action free?
  \item What is a principal bundle map (same group; with $\rho$)? What does the lemma say and how does the proof use freeness and equivariance?
  \item Define the trivial principal bundle and state the section theorem. Sketch both directions, including the role of $\chi_\sigma$ and the identity $\chi_\sigma(p\ract g)=\chi_\sigma(p)\bullet g$.
  \item Why is $LS^2$ not trivial?
\end{enumerate}
\textbf{Concepts.}
\begin{itemize}
  \item Left and right actions, equivariance (\cref{def:pb-left,def:pb-right,prop:pb-right-from-left,def:pb-equivariant}).
  \item Orbits, orbit space, stabilisers, free actions (\cref{def:pb-orbit,def:pb-orbitspace,def:pb-stabiliser,def:pb-free,prop:pb-orbit-G}).
  \item Principal bundles and the frame bundle (\cref{def:pb-principal,def:pb-frame,prop:pb-frame-principal}).
  \item Principal bundle maps and the isomorphism lemma (\cref{def:pb-pbmap,lem:pb-iso}).
  \item Trivial principal bundles and sections (\cref{def:pb-trivial,thm:pb-trivial-section,ex:pb-S2}).
\end{itemize}
\end{keyconcepts}

\flushside
\end{paracol}
